\chapter{Эволюция во времени, сохранение и связанные спины}

Квантовое состояние — не фиксированная точка в гильбертовом пространстве. По мере течения времени оно меняется, но при этом должно меняться так, чтобы не терялись нормировка, вероятности и различимость, закодированная во внутренних произведениях. Эти требования почти полностью определяют форму квантовой динамики. Линейность задаёт оператор временной эволюции, сохранение вероятности делает этот оператор унитарным, а бесконечно малое унитарное преобразование порождается эрмитовым оператором. Этот оператор — гамильтониан. Мы выведем уравнение Schr\"odinger из этих утверждений, увидим, как собственные состояния энергии накапливают фазу, а затем превратим формализм в два конкретных расчёта: прецессию одного спина в магнитном поле и осцилляцию между двумя состояниями пары связанных спинов.

\section{Линейная эволюция в гильбертовом пространстве}

Пусть состояние в момент времени \(t\) равно \(\ket{\psi(t)}\). Мы предполагаем, что его связь с состоянием в нулевой момент времени линейна:
\begin{equation}
\ket{\psi(t)}=U(t)\ket{\psi(0)}.
\end{equation}
Линейность — это постулат, а не теорема, выведенная из чего-то более элементарного. Это экспериментально подтверждённая часть квантовой механики, и попытки заменить её обычно ставят под угрозу либо вероятность, либо законы сохранения. В выбранном базисе \(U(t)\) — просто зависящая от времени матрица.

Двойственный вектор подчиняется сопряжённому правилу. Если
\begin{equation}
\ket{\phi(t)}=U(t)\ket{\phi(0)},
\end{equation}
то
\begin{equation}
\bra{\phi(t)}=\bra{\phi(0)}U^\dagger(t).
\end{equation}
Для матрицы операция эрмитова сопряжения означает комплексное сопряжение с последующим транспонированием. Поэтому бра-вектор подчиняется тому же закону эволюции, что и кет-вектор, только записанному в двойственном пространстве.

\subsection{Сохранение вероятности влечёт унитарность}

Нормированное состояние соответствует полной вероятности, равной единице:
\begin{equation}
\braket{\psi(t)}{\psi(t)}=1.
\end{equation}
Временная эволюция должна сохранять эту норму. Более того, линейное отображение, сохраняющее норму каждого вектора, сохраняет и любое внутреннее произведение. Геометрически, если сохраняются длины двух векторов и их разности, то сохраняется образованный ими треугольник, а значит, и угол между ними. В комплексном векторном пространстве тот же вывод следует, если дополнительно рассмотреть комбинации с множителем \(i\).

Следовательно, для произвольных \(\ket{\phi}\) и \(\ket{\psi}\),
\begin{equation}
\braket{\phi(t)}{\psi(t)}
=
\braket{\phi(0)}{\psi(0)}.
\end{equation}
Подставляя законы эволюции, получаем
\begin{equation}
\bra{\phi(0)}U^\dagger(t)U(t)\ket{\psi(0)}
=
\braket{\phi(0)}{\psi(0)}.
\end{equation}
Поскольку это уравнение выполняется для любой пары векторов,
\begin{equation}
\boxed{U^\dagger(t)U(t)=I}.
\end{equation}
Такой оператор унитарен. Скалярный аналог — комплексное число единичного модуля,
\begin{equation}
e^{i\theta}\bigl(e^{i\theta}\bigr)^*=1.
\end{equation}
Унитарный оператор — матричное обобщение чистой фазы: он может поворачивать состояние в гильбертовом пространстве, но не может растягивать или сжимать его.

\section{Бесконечно малый генератор}

Конечная эволюция складывается из очень малых шагов по времени. Пусть прошедшее время равно \(\epsilon\), настолько малому, что членами порядка \(\epsilon^2\) можно пренебречь. При нулевом прошедшем времени ничего не происходит, так что
\begin{equation}
U(0)=I.
\end{equation}
Поправка первого порядка должна быть пропорциональна \(\epsilon\). Мы записываем её в стандартной форме
\begin{equation}
\boxed{U(\epsilon)=I-\frac{i\epsilon}{\hbar}H}.
\end{equation}
На этом этапе обязателен только множитель \(\epsilon\). \(i\), \(\hbar\) и обозначение \(H\) выбираются так, чтобы объект, остающийся в конце, имел стандартную нормировку и был эрмитовым.

\begin{classroomqa}
\lecturetimestamp{00:12:46}
\audiencequestion{Почему постоянная Планка в привычных единицах оказывается столь поразительно малым числом?}
\lecturerresponse{Численное значение отражает единицы, выбранные людьми. Химики выбрали макроскопические количества, помещающиеся в сосуды, которые можно держать в руке, тогда как атомы несравнимо меньше людей. В единицах, подогнанных под эти макроскопические масштабы, атомный квант действия неизбежно выглядит крошечным. Это не отдельная динамическая загадка; в большой степени это просто исторический и человеческий масштаб, встроенный в систему единиц.}
\end{classroomqa}

Взяв эрмитово сопряжение,
\begin{equation}
U^\dagger(\epsilon)
=I+\frac{i\epsilon}{\hbar}H^\dagger.
\end{equation}
Теперь перемножим и сохраним только члены первого порядка:
\begin{align}
U^\dagger(\epsilon)U(\epsilon)
&=
\left(I+\frac{i\epsilon}{\hbar}H^\dagger\right)
\left(I-\frac{i\epsilon}{\hbar}H\right) \\
&=
I+\frac{i\epsilon}{\hbar}\left(H^\dagger-H\right)
+O(\epsilon^2).
\end{align}
Унитарность требует, чтобы это произведение было равно \(I\), и потому
\begin{equation}
\boxed{H^\dagger=H}.
\end{equation}
Бесконечно малый генератор временной эволюции эрмитов. Поскольку эрмитовы операторы представляют наблюдаемые, любая динамическая система обладает наблюдаемой, которая порождает её эволюцию во времени. Это гамильтониан. Для оператора мы используем \(H\), а \(E\) оставляем для его возможных измеряемых значений, собственных значений.

\begin{classroomqa}
\lecturetimestamp{00:18:34}
\audiencequestion{Почему этот генератор следует называть энергией? Разве энергия — не \(p^2/2m\) или способность совершать работу?}
\lecturerresponse{Это полезные описания в специальных классических ситуациях, а не общее определение. Выражение \(p^2/2m\) есть лишь кинетическая энергия нерелятивистской свободной частицы и не определяет такие вещи, как энергия, запасённая в электромагнитном поле. Для абстрактной системы первое универсальное свойство, которого следует требовать от энергии, — это сохранение. Мы вскоре докажем, что среднее значение \(H\) сохраняется и что фазы его собственных состояний тоже удовлетворяют \(E=\hbar\omega\). Эти два факта и отождествляют генератор с энергией.}
\end{classroomqa}

\section{От малого шага к уравнению Schr\"odinger}

Бесконечно малое изменение уже является дифференциальным уравнением в скрытом виде. При произвольном выборе начала отсчёта времени,
\begin{equation}
\frac{d}{dt}\ket{\psi(t)}
=
\lim_{\epsilon\to0}
\frac{\ket{\psi(t+\epsilon)}-\ket{\psi(t)}}{\epsilon}.
\end{equation}
На доске это сначала записывается при \(t=0\); сдвиг начала отсчёта даёт приведённую выше общую форму.

\begin{figure}[htbp]
\centering
\includegraphics[width=0.72\textwidth]{lecture_09_figure_03.png}
\caption{Разностное отношение, превращающее бесконечно малую эволюцию в дифференциальное уравнение.}
\lectureframe{00:22:15}
\end{figure}

Используем
\begin{equation}
\ket{\psi(t+\epsilon)}
=
\left(I-\frac{i\epsilon}{\hbar}H\right)\ket{\psi(t)}.
\end{equation}
Тогда
\begin{align}
\frac{d}{dt}\ket{\psi(t)}
&=
\lim_{\epsilon\to0}
\frac{\left(I-\frac{i\epsilon}{\hbar}H\right)\ket{\psi(t)}
-\ket{\psi(t)}}{\epsilon} \\
&=
-\frac{i}{\hbar}H\ket{\psi(t)}.
\end{align}
Эквивалентно,
\begin{equation}
\boxed{
i\hbar\frac{d}{dt}\ket{\psi(t)}
=H\ket{\psi(t)}
}.
\end{equation}
Это уравнение Schr\"odinger. Как только известны \(H\) и состояние в один момент времени, оно определяет состояние в более поздние моменты. Вероятности для более поздних измерений затем следуют из обычных правил измерения.

\section{Собственные состояния энергии и вращающаяся фаза}

Самое простое решение начинается с собственного вектора гамильтониана:
\begin{equation}
H\ket{\psi_E}=E\ket{\psi_E}.
\end{equation}

\begin{figure}[htbp]
\centering
\includegraphics[width=0.70\textwidth]{lecture_09_figure_04.png}
\caption{Уравнение на собственные значения энергии на доске.}
\lectureframe{00:28:22}
\end{figure}

Поскольку действие \(H\) на \(\ket{\psi_E}\) оставляет его в том же одномерном направлении, временная эволюция не может смешивать его с другим собственным вектором. Она может только умножать его на скаляр:
\begin{equation}
\ket{\psi_E(t)}=f(t)\ket{\psi_E}.
\end{equation}
Подстановка в уравнение Schr\"odinger даёт
\begin{align}
i\hbar\,\dot f(t)\ket{\psi_E}
&=
f(t)H\ket{\psi_E} \\
&=
E f(t)\ket{\psi_E}.
\end{align}
Следовательно,
\begin{equation}
\dot f(t)=-\frac{iE}{\hbar}f(t),
\qquad
f(t)=e^{-iEt/\hbar},
\end{equation}
и
\begin{equation}
\boxed{
\ket{\psi_E(t)}
=e^{-iEt/\hbar}\ket{\psi_E(0)}
}.
\end{equation}
Собственное состояние энергии стационарно в том смысле, что его луч в гильбертовом пространстве не движется. Тем не менее его кет-вектор накапливает фазу.

Экспоненту можно записать как
\begin{equation}
e^{-iEt/\hbar}
=
\cos\!\left(\frac{Et}{\hbar}\right)
-i\sin\!\left(\frac{Et}{\hbar}\right).
\end{equation}
Знак минус и отсутствие лишнего \(i\) внутри синуса были исправлены в ходе вычисления на доске. Угловая частота равна
\begin{equation}
\omega=\frac{E}{\hbar},
\qquad
\boxed{E=\hbar\omega}.
\end{equation}
Эквивалентно, используя обычную частоту \(\nu=\omega/(2\pi)\), имеем \(E=h\nu\).

\begin{figure}[htbp]
\centering
\includegraphics[width=0.78\textwidth]{lecture_09_figure_05.png}
\caption{Фазовое решение и соотношение энергия--частота выводятся на соседних досках.}
\lectureframe{00:36:19}
\end{figure}

\begin{classroomqa}
\lecturetimestamp{00:34:43}
\audiencequestion{Почему множители \(i\) и \(\hbar\) были введены в бесконечно малый оператор ещё до того, как вычисление их потребовало?}
\lecturerresponse{Это соглашения, выбранные с учётом ожидаемого результата. Без \(i\) унитарность сделала бы генератор антиэрмитовым. Умножение на \(i\) превращает его в эрмитову наблюдаемую с действительными собственными значениями. Множитель \(\hbar\) придаёт этим собственным значениям размерность и нормировочное соглашение энергии. Физическое утверждение — унитарность; эти множители приводят соответствующий генератор к стандартной форме.}
\end{classroomqa}

\begin{classroomqa}
\lecturetimestamp{00:36:18}
\audiencequestion{Относится ли \(E=\hbar\omega\) только к электромагнитным волнам и фотонам?}
\lecturerresponse{Нет. Фотоны — один важный частный случай, но это соотношение более общее. Любое собственное состояние гамильтониана приобретает фазу с угловой частотой \(E/\hbar\). Чтобы связать частоту моды электромагнитного поля с частотой её фотона, требуется дополнительная физика, но само соотношение энергия--фаза уже следует из квантовой временной эволюции.}
\end{classroomqa}

\section{Суперпозиция в энергетическом базисе}

Одно собственное состояние энергии изменяется только на общий фазовый множитель, а общий фазовый множитель не может влиять на вероятность. Интересная динамика начинается, когда состояние содержит несколько собственных векторов энергии. Поскольку эрмитов гамильтониан имеет полный ортонормированный базис собственных векторов,
\begin{equation}
H\ket{\psi_n}=E_n\ket{\psi_n},
\qquad
\ket{\psi(0)}=\sum_n\alpha_n\ket{\psi_n}.
\end{equation}

\begin{figure}[htbp]
\centering
\includegraphics[width=0.72\textwidth]{lecture_09_figure_06.png}
\caption{Разложение произвольного состояния по собственным векторам гамильтониана.}
\lectureframe{00:38:35}
\end{figure}

Линейность позволяет каждой компоненте эволюционировать отдельно:
\begin{equation}
\boxed{
\ket{\psi(t)}
=
\sum_n\alpha_n e^{-iE_n t/\hbar}\ket{\psi_n}
}.
\end{equation}
Компоненты с высокой энергией быстро вращаются на комплексной плоскости; компоненты с низкой энергией вращаются медленнее. На измеримые вероятности влияет не общая для всего состояния фаза, а изменяющиеся относительные фазы между компонентами с разными \(E_n\).

В энергетическом базисе,
\begin{equation}
\begin{aligned}
H&=
\begin{pmatrix}
E_1&0&\cdots\\
0&E_2&\cdots\\
\vdots&\vdots&\ddots
\end{pmatrix},\\[0.6em]
U(t)&=
\begin{pmatrix}
e^{-iE_1t/\hbar}&0&\cdots\\
0&e^{-iE_2t/\hbar}&\cdots\\
\vdots&\vdots&\ddots
\end{pmatrix}.
\end{aligned}
\end{equation}
То же утверждение без выбора базиса кратко записывается как
\begin{equation}
\boxed{U(t)=e^{-iHt/\hbar}},
\end{equation}
для не зависящего от времени гамильтониана.

\begin{classroomqa}
\lecturetimestamp{00:40:17}
\audiencequestion{Является ли это разложение по существу рядом Фурье, и аналогично ли оно разложениям, используемым в теплопроводности?}
\lecturerresponse{Это действительно фурье-подобное разложение по времени: каждому собственному вектору энергии соответствует определённая частота. Методы Фурье и Лапласа также разделяют моды во многих классических уравнениях, включая диффузию. Математическое сходство здесь реально, но утверждения о том, какая именно физическая величина соответствует частоте, нужно проверять в каждом отдельном случае. В квантовой механике конкретное соотношение задаётся уравнением Schr\"odinger.}
\end{classroomqa}

\subsection{Объединение систем}

Для двух невзаимодействующих систем \(A\) и \(B\) суммарный гамильтониан равен сумме отдельных генераторов на тензорном произведении пространств:
\begin{equation}
H_{AB}
=
H_A\otimes I_B+I_A\otimes H_B.
\end{equation}
Если
\begin{equation}
H_A\ket{a}=E_a\ket{a},
\qquad
H_B\ket{b}=E_b\ket{b},
\end{equation}
то
\begin{equation}
H_{AB}\ket{a,b}=(E_a+E_b)\ket{a,b}.
\end{equation}
Энергии, а следовательно, и фазовые частоты, для разделённых невзаимодействующих систем складываются. Взаимодействие даёт дополнительный член \(H_{\mathrm{int}}\), так что эта простая сумма уже не даёт полной энергии.

\begin{classroomqa}
\lecturetimestamp{00:56:20}
\audiencequestion{Если две отдельные системы содержат компоненты с одной и той же частотой, что происходит, когда эти системы объединяют?}
\lecturerresponse{Для невзаимодействующих систем собственные значения энергии объединённой системы равны суммам отдельных собственных значений. Совпадение двух частот может привести к вырождению, но само по себе не смешивает соответствующие собственные векторы. Для смешивания нужен член взаимодействия или выбор базиса внутри вырожденного подпространства. Как только системы начинают взаимодействовать, расширенный гамильтониан нужно диагонализовать как новую задачу.}
\end{classroomqa}

\section{Средние значения и коммутаторы}

Мы прослеживали эволюцию состояния, считая наблюдаемые фиксированными. Это картина Schr\"odinger. Пусть \(A\) — наблюдаемая, не имеющая явной зависимости от времени. Её среднее значение имеет вид
\begin{equation}
\langle A\rangle_t
=
\bra{\psi(t)}A\ket{\psi(t)}.
\end{equation}
Продифференцируем как бра-вектор, так и кет-вектор:
\begin{equation}
\frac{d}{dt}\langle A\rangle_t
=
\bra{\dot\psi}A\ket{\psi}
+
\bra{\psi}A\ket{\dot\psi}.
\end{equation}
Уравнение Schr\"odinger и эрмитово-сопряжённое к нему дают
\begin{equation}
\ket{\dot\psi}
=-\frac{i}{\hbar}H\ket{\psi},
\qquad
\bra{\dot\psi}
=\frac{i}{\hbar}\bra{\psi}H.
\end{equation}
Следовательно,
\begin{align}
\frac{d}{dt}\langle A\rangle_t
&=
\frac{i}{\hbar}\bra{\psi}HA\ket{\psi}
-\frac{i}{\hbar}\bra{\psi}AH\ket{\psi} \\
&=
\frac{i}{\hbar}\bra{\psi}(HA-AH)\ket{\psi}.
\end{align}
Определим коммутатор
\begin{equation}
[H,A]=HA-AH.
\end{equation}
Тогда уравнение движения для среднего значения имеет вид
\begin{equation}
\boxed{
\frac{d}{dt}\langle A\rangle_t
=
\frac{i}{\hbar}\langle[H,A]\rangle_t
}.
\end{equation}

\begin{figure}[htbp]
\centering
\includegraphics[width=0.74\textwidth]{lecture_09_figure_07.png}
\caption{Уравнение для среднего значения и коммутатор \(AB-BA=[A,B]\).}
\lectureframe{00:51:05}
\end{figure}

\begin{classroomqa}
\lecturetimestamp{00:49:13}
\audiencequestion{Почему коммутатор в этом уравнении сам стоит под знаком среднего значения?}
\lecturerresponse{Коммутатор — это оператор. Любой оператор, помещённый между \(\bra{\psi}\) и \(\ket{\psi}\), имеет среднее значение, так что \(\bra{\psi}[H,A]\ket{\psi}\) — это просто среднее этого коммутатора в текущем состоянии.}
\end{classroomqa}

\subsection{Законы сохранения}

Положим \(A=H\). Оператор всегда коммутирует сам с собой:
\begin{equation}
[H,H]=HH-HH=0.
\end{equation}
Следовательно,
\begin{equation}
\boxed{
\frac{d}{dt}\langle H\rangle_t=0
}.
\end{equation}
Среднее значение гамильтониана сохраняется. В этом и состоит обещанная структурная причина отождествления \(H\) с энергией. У всякой изолированной системы есть гамильтониан, и любой не зависящий от времени гамильтониан даёт этот закон сохранения.

В более общем случае,
\begin{equation}
[H,A]=0
\quad\Longrightarrow\quad
\frac{d}{dt}\langle A\rangle_t=0.
\end{equation}

\begin{classroomqa}
\lecturetimestamp{00:51:41}
\audiencequestion{Каков физический смысл совокупности всех операторов, коммутирующих с \(H\), и включает ли она величины помимо энергии?}
\lecturerresponse{Это сохраняющиеся наблюдаемые системы. Какие именно из них существуют, зависит от гамильтониана. Свободная частица сохраняет импульс, тогда как сила может нарушить это сохранение. Энергия присутствует всегда, потому что \(H\) коммутирует сам с собой; дополнительные коммутирующие операторы выражают дополнительные симметрии и законы сохранения.}
\end{classroomqa}

\section{Один спин в магнитном поле}

Рассмотрим спин с операторами Паули
\begin{equation}
\boldsymbol{\sigma}
=(\sigma_1,\sigma_2,\sigma_3)
\end{equation}
в фиксированном внешнем магнитном поле \(\mathbf B\). Магнитный момент пропорционален спину. В используемой здесь конвенции,
\begin{equation}
\boxed{
H=\frac{\mu}{2}\boldsymbol{\sigma}\cdot\mathbf B
}.
\end{equation}
Множитель \(1/2\), знак и зависящий от заряда коэффициент можно включить в определение \(\mu\). Поле рассматривается классически: большой устойчивый магнит создаёт \(\mathbf B\), и квантовое состояние этого магнита мы не моделируем.

\begin{figure}[htbp]
\centering
\includegraphics[width=0.68\textwidth]{lecture_09_figure_08.png}
\caption{Оператор магнитного момента взаимодействует с внешним магнитным полем.}
\lectureframe{01:01:45}
\end{figure}

Выберем поле направленным вдоль третьей оси,
\begin{equation}
\mathbf B=B\hat{\mathbf z},
\qquad
H=\frac{\mu B}{2}\sigma_3.
\end{equation}
Будем отслеживать средние значения трёх компонент спина. Чтобы сократить запись уравнений, введём
\begin{equation}
s_i(t)=\langle\sigma_i\rangle_t.
\end{equation}
Для продольной компоненты имеем
\begin{equation}
\dot s_3
=
\frac{i}{\hbar}\langle[H,\sigma_3]\rangle
=0,
\end{equation}
поскольку \(H\) пропорционален \(\sigma_3\). Проекция спина на направление поля фиксирована. Геометрически конец вектора среднего спина должен оставаться на конусе с фиксированным углом раскрытия вокруг оси поля.

\begin{figure}[htbp]
\centering
\includegraphics[width=0.58\textwidth]{lecture_09_figure_09.png}
\caption{При фиксированной продольной компоненте спина движение остаётся на конусе вокруг оси магнитного поля.}
\lectureframe{01:04:45}
\end{figure}

Для первой поперечной компоненты,
\begin{equation}
\dot s_1
=
\frac{i\mu B}{2\hbar}
\langle[\sigma_3,\sigma_1]\rangle.
\end{equation}
Коммутатор можно непосредственно проверить, используя
\begin{equation}
\begin{aligned}
\sigma_3&=
\begin{pmatrix}1&0\\0&-1\end{pmatrix},
&
\sigma_1&=
\begin{pmatrix}0&1\\1&0\end{pmatrix},\\[0.5em]
\sigma_2&=
\begin{pmatrix}0&-i\\i&0\end{pmatrix}.
\end{aligned}
\end{equation}
Действительно,
\begin{equation}
\sigma_3\sigma_1=i\sigma_2,
\qquad
\sigma_1\sigma_3=-i\sigma_2,
\end{equation}
поэтому
\begin{equation}
[\sigma_3,\sigma_1]=2i\sigma_2.
\end{equation}
Следовательно,
\begin{equation}
\dot s_1=-\frac{\mu B}{\hbar}s_2.
\end{equation}
Аналогично, используя \([\sigma_3,\sigma_2]=-2i\sigma_1\),
\begin{equation}
\dot s_2=\frac{\mu B}{\hbar}s_1.
\end{equation}

\begin{figure}[htbp]
\centering
\includegraphics[width=0.72\textwidth]{lecture_09_figure_10.png}
\caption{Вычисление коммутатора матриц Паули, приводящее к поперечному уравнению движения.}
\lectureframe{01:09:35}
\end{figure}

Положим
\begin{equation}
\Omega=\frac{\mu B}{\hbar}.
\end{equation}
Три уравнения имеют вид
\begin{equation}
\dot s_1=-\Omega s_2,
\qquad
\dot s_2=\Omega s_1,
\qquad
\dot s_3=0.
\end{equation}
Решение с поперечной амплитудой \(s_\perp\) таково:
\begin{equation}
s_1(t)=s_\perp\cos(\Omega t+\varphi),
\qquad
s_2(t)=s_\perp\sin(\Omega t+\varphi),
\end{equation}
при этом \(s_3\) постоянно.

\begin{figure}[htbp]
\centering
\includegraphics[width=0.64\textwidth]{lecture_09_figure_11.png}
\caption{Поперечные компоненты спина образуют синусоидальную--косинусоидальную пару с угловой частотой \(\mu B/\hbar\).}
\lectureframe{01:11:45}
\end{figure}

Следовательно, средний спин прецессирует вокруг магнитного поля, как гироскоп. Выбор поля вдоль \(z\) был лишь удобством. Для поля в любом направлении компонента, параллельная \(\mathbf B\), сохраняется, а перпендикулярная вращается вокруг этого направления.

\begin{classroomqa}
\lecturetimestamp{01:13:02}
\audiencequestion{Не должно ли решение в виде синуса и косинуса содержать дополнительную фазу?}
\lecturerresponse{Да. Фаза \(\varphi\) задаёт начальное направление в поперечной плоскости или, что эквивалентно, выбор момента, обозначенного \(t=0\). Приравнивание её к нулю просто означает выбор начального положения спина вдоль одной конкретной поперечной оси.}
\end{classroomqa}

\subsection{Подготовка и измерение}

Предположим, что ансамбль подготовлен в собственном состоянии \(+1\) оператора \(\sigma_1\). При \(t=0\) измерение \(\sigma_1\) заведомо даёт \(+1\), тогда как измерения \(\sigma_2\) или \(\sigma_3\) дают равные вероятности для \(\pm1\). Через четверть периода прецессии состояние повернётся в собственное состояние \(+1\) оператора \(\sigma_2\): тогда результат измерения \(\sigma_2\) определён, а для \(\sigma_1\) неопределённость максимальна. Ещё через четверть периода оно достигает собственного состояния \(-1\) оператора \(\sigma_1\). Поэтому вращающиеся средние значения имеют прямой операциональный смысл.

\begin{classroomqa}
\lecturetimestamp{01:14:04}
\audiencequestion{Если для средних значений нужно много электронов, не будут ли члены ансамбля иметь разные фазы и тем самым смазывать прецессию?}
\lecturerresponse{Нет, если над каждым членом проводится один и тот же эксперимент. Одинаковая подготовка даёт каждому электрону одно и то же начальное состояние и, следовательно, одну и ту же фазовую эволюцию. Если их начальные фазы не контролируются, ансамбль не представляет собой чистую подготовку; это смешанное состояние, и дефазировка может уменьшить поперечное среднее.}
\end{classroomqa}

\section{Два связанных спина}

Теперь рассмотрим два спина. Одного лишь утверждения, что они запутаны, недостаточно, чтобы предсказать, как они будут себя вести. Для задания динамики нужен гамильтониан.

\begin{classroomqa}
\lecturetimestamp{01:16:49}
\audiencequestion{Если два электрона запутаны так, что вдоль любой общей оси получаются противоположные результаты, значит ли это, что их движения всегда сдвинуты по фазе на \(180^\circ\)?}
\lecturerresponse{Из одной только запутанности никакого динамического вывода не следует. Указанная антикорреляция описывает синглетное состояние, но его последующая эволюция зависит от оператора энергии. Сначала нужно указать, имеется ли внешнее поле, взаимодействие между спинами или и то и другое.}
\end{classroomqa}

Рассмотрим простое изотропное взаимодействие, в котором каждый спин создаёт эффективное поле, испытываемое другим:
\begin{equation}
\boxed{
H=J\,\boldsymbol{\sigma}\cdot\boldsymbol{\tau}
=
J\left(
\sigma_1\tau_1+\sigma_2\tau_2+\sigma_3\tau_3
\right)
}.
\end{equation}
Операторы \(\boldsymbol{\sigma}\) действуют на первый спин, а \(\boldsymbol{\tau}\) — на второй. Константа связи \(J\) задаёт энергетический масштаб; её величина и знак здесь не уточняются.

\begin{figure}[htbp]
\centering
\includegraphics[width=0.75\textwidth]{lecture_09_figure_12.png}
\caption{Изотропное взаимодействие двух спинов \(H\propto\boldsymbol{\sigma}\cdot\boldsymbol{\tau}\).}
\lectureframe{01:19:20}
\end{figure}

\begin{classroomqa}
\lecturetimestamp{01:17:49}
\audiencequestion{Обязательно ли при помещении двух запутанных электронов в однородное магнитное поле их запутанность исчезнет?}
\lecturerresponse{Не обязательно. Поле изменяет гамильтониан и может стремиться выровнять спины, особенно если оно сильное, но унитарная эволюция в однородном поле сама по себе не гарантирует исчезновения запутанности. Ответ зависит от полного гамильтониана и начального состояния.}
\end{classroomqa}

\subsection{Энергии триплета и синглета}

Пространство Гильберта двух спинов имеет размерность 4. В триплет--синглетном базисе взаимодействие диагонально:
\begin{align}
\ket{T_+}&=\ket{\uparrow\uparrow}, \\
\ket{T_0}&=
\frac{\ket{\uparrow\downarrow}+\ket{\downarrow\uparrow}}{\sqrt2}, \\
\ket{T_-}&=\ket{\downarrow\downarrow}, \\
\ket{S}&=
\frac{\ket{\uparrow\downarrow}-\ket{\downarrow\uparrow}}{\sqrt2}.
\end{align}
Три триплетных состояния имеют одну и ту же энергию \(E_T\); синглет имеет другую энергию \(E_S\). Следовательно, существуют четыре независимых собственных вектора, но только два различных уровня энергии. Триплетный уровень трёхкратно вырожден.

\begin{figure}[htbp]
\centering
\includegraphics[width=0.74\textwidth]{lecture_09_figure_13.png}
\caption{Гамильтониан связанных спинов имеет вырожденный триплетный уровень и отдельный синглетный уровень.}
\lectureframe{01:20:30}
\end{figure}

\begin{classroomqa}
\lecturetimestamp{01:21:16}
\audiencequestion{Что будет, если начальное состояние уже является синглетным собственным состоянием энергии?}
\lecturerresponse{Оно остаётся синглетом и приобретает лишь фазовый множитель \(e^{-iE_St/\hbar}\). Как и у любого отдельного собственного состояния энергии, его измеримые свойства не меняются просто потому, что весь кет-вектор умножается на этот общий фазовый множитель.}
\end{classroomqa}

\begin{classroomqa}
\lecturetimestamp{01:21:54}
\audiencequestion{Во время этого расчёта сообщает ли внешнее магнитное поле энергию двухспиновой системе?}
\lecturerresponse{В решаемом примере внешнего поля нет. Единственное поле, о котором здесь идёт речь, — это взаимное взаимодействие, представляемое \(J\,\boldsymbol{\sigma}\cdot\boldsymbol{\tau}\). Добавление внешнего поля добавило бы ещё один член в \(H\) и означало бы уже другую задачу.}
\end{classroomqa}

\subsection{Начнём с одного спина вверх и одного вниз}

Собственное состояние энергии лишь набирало бы фазу, поэтому выберем продуктное состояние, которое не является собственным состоянием энергии:
\begin{equation}
\ket{\psi(0)}=\ket{\uparrow\downarrow}.
\end{equation}
Разложим его по двум соответствующим собственным векторам:
\begin{equation}
\boxed{
\ket{\uparrow\downarrow}
=
\frac{1}{\sqrt2}\ket{T_0}
+
\frac{1}{\sqrt2}\ket{S}
}.
\end{equation}

\begin{figure}[htbp]
\centering
\includegraphics[width=0.76\textwidth]{lecture_09_figure_14.png}
\caption{Продуктное состояние \(\ket{\uparrow\downarrow}\) разложено на триплетную и синглетную компоненты.}
\lectureframe{01:24:30}
\end{figure}

Теперь каждая энергетическая компонента получает свой фазовый множитель:
\begin{equation}
\ket{\psi(t)}
=
\frac{e^{-iE_Tt/\hbar}}{\sqrt2}\ket{T_0}
+
\frac{e^{-iE_St/\hbar}}{\sqrt2}\ket{S}.
\end{equation}

\begin{classroomqa}
\lecturetimestamp{01:25:59}
\audiencequestion{Какая энергия больше: триплетная энергия \(E_T\) или синглетная энергия \(E_S\)?}
\lecturerresponse{Это зависит от знака константы связи. Для \(H=J\,\boldsymbol{\sigma}\cdot\boldsymbol{\tau}\) при \(J>0\) синглет ниже, а триплет выше; при смене знака порядок меняется на обратный. Добавление константы к обеим энергиям меняет только общую фазу, так что любой из уровней можно положить равным нулю, не изменяя вероятностей перехода.}
\end{classroomqa}

Вернёмся к продуктному базису:
\begin{align}
\ket{\psi(t)}
&=
\frac12
\left(
e^{-iE_Tt/\hbar}+e^{-iE_St/\hbar}
\right)\ket{\uparrow\downarrow} \\
&\quad+
\frac12
\left(
e^{-iE_Tt/\hbar}-e^{-iE_St/\hbar}
\right)\ket{\downarrow\uparrow}.
\end{align}

\begin{figure}[htbp]
\centering
\includegraphics[width=0.76\textwidth]{lecture_09_figure_15.png}
\caption{Эволюционировавшее состояние, вновь собранное в базисе \(\ket{\uparrow\downarrow}\), \(\ket{\downarrow\uparrow}\).}
\lectureframe{01:30:40}
\end{figure}

При \(t=0\) коэффициент при \(\ket{\downarrow\uparrow}\) обращается в нуль, как и должно быть. В более поздние моменты времени относительная фаза между синглетной и триплетной компонентами приводит к росту этого коэффициента. Введём среднюю энергию и расщепление,
\begin{equation}
\overline E=\frac{E_T+E_S}{2},
\qquad
\Delta E=E_T-E_S.
\end{equation}
Вынося фазу, соответствующую средней энергии, получаем
\begin{equation}
\boxed{
\begin{aligned}
\ket{\psi(t)}
&=e^{-i\overline E t/\hbar}\Biggl[
\cos\!\left(\frac{\Delta E\,t}{2\hbar}\right)
\ket{\uparrow\downarrow}
\\
&\qquad
-i\sin\!\left(\frac{\Delta E\,t}{2\hbar}\right)
\ket{\downarrow\uparrow}
\Biggr]
\end{aligned}
}.
\end{equation}

\begin{figure}[H]
\centering
\includegraphics[width=0.76\textwidth]{lecture_09_figure_16.png}
\caption{Средняя энергия становится общей фазой, тогда как синглет--триплетное расщепление проявляется в амплитудах синуса и косинуса.}
\lectureframe{01:35:15}
\end{figure}

Множитель \(e^{-i\overline E t/\hbar}\) общий для всего кет-вектора и сокращается во всех вероятностях. Измеримые вероятности равны
\begin{equation}
P_{\uparrow\downarrow}(t)
=
\cos^2\!\left(\frac{\Delta E\,t}{2\hbar}\right),
\qquad
P_{\downarrow\uparrow}(t)
=
\sin^2\!\left(\frac{\Delta E\,t}{2\hbar}\right).
\end{equation}
Состояние осциллирует от \(\ket{\uparrow\downarrow}\) к \(\ket{\downarrow\uparrow}\) и обратно с частотой, определяемой только разностью энергий. В промежуточные моменты времени обе амплитуды ненулевые, и, за исключением особых моментов, два спина запутаны.

Этот расчёт иллюстрирует общий метод. Выберите базис, в котором поставлен физический вопрос, разложите начальное состояние по собственным векторам гамильтониана, дайте каждому из них приобрести свою фазу, а затем вновь соберите части в базисе измерения. После диагонализации гамильтониана квантовая динамика становится систематической; всё наблюдаемое изменение обусловлено относительной фазой.
